% !Mode:: "TeX:UTF-8"

\chapter{前后端解耦架构}

\section{前后端解耦架构设计}

本文实现的乱序处理器采用前后端解耦架构。前端（FrontEnd）负责取指与分支预测，后端（BackEnd）负责译码、寄存器重命名、分派、发射、执行、写回、晚执行、提交。前后端之间通过一个指令队列（Instruction Queue）进行缓冲与解耦，后端在需要冲刷流水线时通过 flush 信号同步清空该队列，分支预测失败或异常恢复时通过 redirect 信号将正确的 PC 反馈给前端。

\begin{figure}[H]
  \centering
  \includegraphics{fe-be-arch.png}
  \caption{前后端解耦架构与指令队列}
  \label{fig:fe-be-arch}
\end{figure}

如图~\ref{fig:fe-be-arch} 所示，前端产生的指令流经指令队列缓冲后进入后端；后端在检测到分支预测错误或异常时，通过 flush 信号清空队列中的无效指令，同时通过 redirect 信号将恢复后的目标 PC 发送给前端，使前端从正确地址重新开始取指。

\section{前后端解耦架构的含义}

前后端解耦的含义体现在时序解耦与实现解耦两个层面，并由此带来弹性的设计空间。

在时序解耦方面，前后端之间的连接方式是可配置的：既可以选择点对点直连（组合逻辑穿透，零周期延迟），也可以插入深度可调的缓冲队列。当使用队列时，前后端各自按自身节奏运行——后端因数据依赖或功能单元冲突暂停消费时，前端仍可继续取指填充队列，避免取指带宽浪费；前端遭遇 ICache 缺失时，后端可继续执行队列中已缓冲的指令，减少流水线气泡。队列深度作为参数，可根据实际时序需求灵活调整。

在实现解耦方面，前端与后端通过标准 ready-valid 握手协议通信，两侧的内部实现完全独立。前端既可以采用流水化结构，也可以完全使用软件实现。在开发过程中，硬件前端由于需要等待总线响应与 ICache 缺失处理，指令供应速率有限，后端负载偏低，难以充分暴露乱序流水线在高压力下的时序问题。为此，本课题利用解耦架构的可替换性，在前端内嵌入 Spike 参考模型：Spike 每执行一条指令即可获得精确的下一条指令 PC，使分支预测几乎 100\% 正确，同时取指不经过总线，几乎每个周期都能向后端供应一条有效指令。这一做法相当于对后端进行了满载压力测试，成功暴露了一个在常规硬件前端下因指令供应不足、未形成反压（back-pressure）而被隐藏的 bug——该 bug 表现为反压条件下数据转发时机错过，仅在后端持续满载时才会触发。定位并修复该问题后，前端切换回正常硬件实现。图~\ref{fig:spike-frontend} 展示了使用 Spike 软件前端运行 MicroBench 测试套件的结果：全部测试项通过，共提交 624375 条指令，其中 142230 条分支指令仅产生 514 次预测错误，命中率达 99.64\%，验证了该方案的指令供应充分性。这一经验表明，解耦架构不仅降低了模块间的开发耦合度，还为针对性的压力测试提供了灵活的验证手段。

\begin{figure}[H]
  \centering
  \includegraphics{spike-frontend-bench.png}
  \caption{Spike 软件前端运行 MicroBench 的测试结果（分支预测命中率 99.64\%）}
  \label{fig:spike-frontend}
\end{figure}

\section{本章小结}

弹性设计空间是解耦架构的另一直接收益。指令队列作为缓冲层，使得前端分支预测器或取指流水线的升级不影响后端接口，有利于微架构的增量迭代与独立优化。
